\section{From a Local Attribution Direction to Multi-Round Selection}
\label{sec:single-to-multi}

Equation~\ref{eq:rho-vf-equivalence} is a local statement: at a fixed
$\vomega$, a two-model loss difference supplies a descent direction for a
finite-penalty surrogate. Data selection is instead a constrained, path-
dependent decision problem. It must specify an initial state, how long the
real target is trained before the selector changes, whether the same outer
state is retained, and how a continuous state becomes a discrete subset.

\paragraph{A common warm start.}
All comparisons begin with the same target initialization and the same random
budget-$B$ subset $S_0$. The target trains on $S_0$ for $E_0$ epochs before any
method-dependent selection. This avoids choosing an arbitrary Top-$B$ set from
tied uniform weights and makes the first comparison occur at an identical,
nontrivial target state.

\paragraph{One fixed horizon, one to five selection rounds.}
After the warm start, target training is partitioned into five equal blocks of
$L$ epochs, for a common horizon $H=5L$. At the start of a block, one complete
selection round (i) updates $\thetahat$ and $\thetatilde$ under a fixed
inner-pass rule, (ii) computes a global direction, (iii) updates and projects
the continuous state, and (iv) deploys
$S_{r}=\operatorname{TopB}(\vomega_r)$ for real target training. We compare
$R\in\{1,2,3,4,5\}$: the first $R$ boundaries update the selector, after which
its continuous state and subset are frozen while the target still trains to
$H$. An extra round is therefore a complete bilevel selection update, not an
extra target-SGD step or a virtual epoch that is discarded.

\paragraph{Reset separates rescoring from outer memory.}
A reset branch computes every new direction from
$\vomega^{\mathrm{unif}}$; a persistent branch updates the same
$\vomega_r\in\Wbudget$ across boundaries. If both improve with $R$, repeated
evaluation at a changing learner state is sufficient. If only the persistent
branch improves, the evidence instead favors optimization of continuous
weights along a path. In both cases Top-$B$ is a deployment operator: the
binary mask used to train the target is never used as the next differentiable
outer state.

\paragraph{Where can a multi-round gain come from?}
Relative to one-shot selection, four state changes can matter:
\begin{enumerate}
    \item the target model changes, so the value of an example is conditional
    on a new learner state;
    \item the two inner models track new finite-penalty solutions;
    \item the continuous outer state remembers previous directions; and
    \item Top-$B$ is applied again, changing the deployed discrete subset.
\end{enumerate}
The fixed-horizon reset/persistent comparison identifies the first three at
the method level; score, state, subset, and target checkpoints at every block
boundary expose the fourth. A claim that ``more rounds help'' is incomplete
unless the responsible transition is identified.

\paragraph{RHO as an approximation endpoint.}
Practical RHO removes two costly pieces of the multi-round VF procedure: it
uses the current nonconverged target instead of $\thetahat(\vomega)$ and a
frozen holdout-only IL model instead of
$\thetatilde(\vomega)$. It then scores only a random large batch and
immediately trains on its top fraction. We therefore compare RHO neither as an
unrelated baseline nor as an exact VF implementation, but as the efficient end
of a nested path. Global scoring, local quota, immediate updates, online
rescoring, and comparator choice are switched on separately in
Section~\ref{sec:experiments}.
